\section{Experiments}
\label{sec:experiments}
\paragraph{Models, prompts, and fidelity.}
We adopt the protocol of RA-CFGCache Table~1 \citep{liu2026racfgcache} so that our numbers can be placed next to the three closest training-free baselines without re-tuning them. CogVideoX-2B \citep{yang2025cogvideox} generates 49 frames at $720\times480$ and 8\,fps with 50 sampling steps and guidance scale 6. Wan2.1-T2V-1.3B \citep{wan2025wan} generates 81 frames at $832\times480$ and 16\,fps with 50 UniPC steps \citep{zhao2023unipc}, flow shift 3, and guidance scale 5. Prompts are the non-empty lines of the VBench prompt resource \citep{huang2024vbench} used by that harness, of which there are 100; prompt $k$ uses seed $k$, negative prompt ``animation'', and batch size one. Every accelerated video is compared with the full 50-step video of the same prompt and seed by LPIPS (AlexNet), SSIM, and PSNR. Using the VBench prompt list does not constitute a VBench evaluation; reference-free VBench dimensions are reported separately when they are run.

\paragraph{Baselines.}
TeaCache \citep{liu2025teacache} uses thresholds 0.2 (CogVideoX) and 0.15 (Wan); TaylorSeer \citep{liu2025taylorseer} uses interval 3, order 1, and three warm-up steps; RA-CFGCache-Slow/Fast use thresholds 0.6/1.4 (CogVideoX) and 0.4/0.6 (Wan) with the released calibration tables. All baselines run inside the pinned public RA-CFGCache harness, which also contains the TeaCache and TaylorSeer implementations, with two corrections that do not change any method. First, the pinned sampling script forwarded the command-line threshold under a key that the RA-CFGCache controller does not read, so both nominal variants silently ran at the default threshold; we bind the threshold explicitly and assert the effective value at run time. Second, all denoising and caching run on the GPU without CPU offload.

\paragraph{Hardware deviations.}
The published numbers use an H200 in bf16. We run the identical protocol on four H20-96GB cards in bf16, with no CPU offload, no VAE tiling, and no text-encoder parking, so the precision matches the paper and the only deviation is the accelerator itself. A single V100-32GB in fp16 was used as a cross-check: it reproduces CogVideoX and the four Wan rows at $n{=}32$, but its 31.7\,GiB cannot hold TaylorSeer's per-layer derivative cache on Wan, which needs roughly 90\,GiB across the two CFG branches. Latencies are chunk averages measured while four cards are busy, so they are provisional; the quality columns are the comparable ones. All wall-clock numbers for the controller, including its online-identification overhead, will be re-timed on an otherwise idle device.

\begin{table*}[t]
\centering
\small
\setlength{\tabcolsep}{4.2pt}
\caption{Paired 100-prompt baseline reproduction on CogVideoX-2B and Wan2.1-T2V-1.3B (ours: H20-96GB \texttt{bfloat16}; paper: published H200 \texttt{bfloat16} values from \citet{liu2026racfgcache}). RA-CFGCache-Slow achieves the strongest reconstruction fidelity across both backbones; $^\dagger$speedup ratios reflect different GPU architectures.}
\label{tab:baselines}
\begin{tabular}{llccccc}
\toprule
Model & Method & $n$ & LPIPS$\downarrow$ & SSIM$\uparrow$ & PSNR$\uparrow$ & Speedup$^\dagger$ \\
 & & & ours / paper & ours / paper & ours / paper & ours / paper \\
\midrule
CogVideoX-2B & Full ($T{=}50$) & 100 & 0.0000 / 0.0000 & 1.0000 / 1.0000 & $\infty$ / $\infty$ & 1.00$\times$ / 1.00$\times$ \\
 & Fewer steps ($T{=}25$) & 100 & 0.4704 / 0.4751 & 0.5350 / 0.5445 & 14.08 / 13.61 & 1.89$\times$ / 1.86$\times$ \\
 & TeaCache & 100 & 0.2122 / 0.2040 & 0.7392 / 0.7538 & 21.22 / 21.48 & 2.64$\times$ / 2.42$\times$ \\
 & TaylorSeer ($N{=}3,O{=}1$) & 100 & 0.4047 / -- & 0.5626 / -- & 15.38 / -- & 2.36$\times$ / -- \\
 & RA-CFGCache-Slow & 100 & 0.0990 / 0.0968 & 0.8661 / 0.8746 & 26.89 / 26.93 & 2.63$\times$ / 2.39$\times$ \\
 & RA-CFGCache-Fast & 100 & 0.2092 / 0.2052 & 0.7493 / 0.7614 & 21.88 / 21.89 & 3.74$\times$ / 3.13$\times$ \\
\midrule
Wan2.1-T2V-1.3B & Full ($T{=}50$) & 100 & 0.0000 / 0.0000 & 1.0000 / 1.0000 & $\infty$ / $\infty$ & 1.00$\times$ / 1.00$\times$ \\
 & Fewer steps ($T{=}25$) & 100 & 0.4204 / 0.4371 & 0.5700 / 0.5658 & 15.85 / 15.87 & 1.96$\times$ / 1.93$\times$ \\
 & TeaCache & 100 & 0.1681 / 0.1661 & 0.7887 / 0.7955 & 22.84 / 23.08 & 2.53$\times$ / 2.44$\times$ \\
 & TaylorSeer ($N{=}3,O{=}1$) & 100 & 0.3496 / 0.3632 & 0.6283 / 0.6187 & 17.10 / 16.96 & 2.47$\times$ / 2.13$\times$ \\
 & RA-CFGCache-Slow & 100 & 0.0941 / 0.0993 & 0.8820 / 0.8778 & 28.11 / 28.03 & 2.65$\times$ / 2.56$\times$ \\
 & RA-CFGCache-Fast & 100 & 0.1352 / 0.1409 & 0.8420 / 0.8386 & 26.14 / 26.12 & 3.14$\times$ / 2.99$\times$ \\
\bottomrule
\end{tabular}
\end{table*}


\paragraph{Reproduced baselines.}
Table~\ref{tab:baselines} reports all 100 prompts of both backbones; every row has 100 paired samples and no row is imputed. Across the five rows that the published table also contains, LPIPS differs by at most 0.017, SSIM by at most 0.015 and PSNR by at most 0.5\,dB, and the ordering is identical: RA-CFGCache-Slow is the strongest fidelity baseline on both models, TeaCache and RA-CFGCache-Fast trade places with the backbone, and 25 steps is the weakest non-trivial row. This makes RA-CFGCache-Slow the primary comparison for any matched-budget claim. The speedups are not comparable to the published ones: they are chunk averages on four busy H20 cards, and H20 is not H200. TaylorSeer on Wan is reproduced here for the first time in this project ($0.3496$ LPIPS at $2.47\times$); it is clearly worse than TeaCache at a comparable budget, the same pattern as the published Wan row, so its poor CogVideoX result is a property of the method's first-order time extrapolation and not an implementation error. One prompt (id 56 on CogVideoX) collapses to a white clip for the reference and for every method; excluding it moves the row means by at most 0.004 LPIPS and changes no ordering. No statistical gate is claimed for the baselines.

\paragraph{Development platform and component gates.}
Components of \method\ are first tested on a cheaper development setting: CogVideoX-2B, 17 frames at $720\times480$, 30 DDIM steps \citep{song2021ddim}, guidance 6, and eight fixed prompts with two seeds each (16 paired samples). These prompts are disjoint from the evaluation prompts. Each component must pass a pre-registered gate before it enters the full controller. For the forecaster, all methods compute exactly the same steps, given by the static TeaCache schedule at $N\in\{10,13,17\}$ exact evaluations, and differ only in what skipped steps return. Fidelity is the relative $\ell_2$ error of the terminal latent against the all-exact run. The forecaster passes if no sample diverges and the paired median error decreases by at least 10\% relative to the better of reuse and first-order temporal extrapolation at every budget. For propagation pricing, a causal estimate must rank the terminal effect of single-step perturbations better than a leave-one-prompt-out timestep prior of the RA-CFGCache type. The full controller is then compared at matched exact-evaluation budgets with the same forecasts under fixed and prior-based schedules, with paired tests requiring at least 16 samples, a median improvement of at least 25\%, a win rate of at least 0.75, and an effect larger than twice its standard error.

\paragraph{Limitations.}
\textbf{\method\ has no validated result.} Both pre-registered component gates fail: the online propagation price ranks terminal error at Spearman $0.33$ against $0.95$ for a timestep prior, and the state-conditioned forecaster beats the better of reuse and first-order extrapolation by $9.1\%$ at $N{=}13$ but loses $22\%$ at $N{=}10$. The controller has therefore never been run end-to-end at the scale of Table~\ref{tab:baselines}, which reports baselines only. The reproduction also runs on four H20 cards rather than the published H200, so only the quality columns are comparable; one CogVideoX prompt collapses to a white clip for the reference and for every method; and each baseline is evaluated at the budget its own scheduler selects, so its wall-clock advantage over a fixed budget is not uniform.
